\section{Details of the Benchmark}
\label{sec:benchmark-details}

\subsection{Data source and construction}
The in-distribution (IID) benchmark is built from the source cohort abdominal--pelvic CT
studies with paired radiology reports. Each study contributes one CT volume
(\texttt{.nii.gz}) and a free-text report; the report's organ-section prose is the
source of all labels. A rule-based labeler (\texttt{rule\_v4}) first extracts
candidate labels, which are then verified by an ensemble of LLM readers
(\S\ref{sec:verification}). The released benchmark contains
\textbf{280{,}111 questions over 25{,}270 studies}, partitioned into training,
validation, and test splits. Questions are answered with a single symbol
(\texttt{a}/\texttt{b} for yes/no, a letter for multiple choice, or \texttt{1}--\texttt{5}
for severity), so every item is a closed-form decision rather than free-text
generation.

\subsection{Task families and question types}
The benchmark spans six question families (Table~\ref{tab:families}). The
\emph{finding} family asks whether a specific abnormality is present (binary
yes/no); the \emph{organ} family asks for the most prominent abnormality of a given
organ system as a multiple-choice question; the \emph{score} family grades severity
on an ordinal L0--L4 scale; the \emph{side} family asks laterality; the
\emph{triage} family asks whether the impression indicates an acute intra-abdominal
process; and an \emph{extended} family adds report-mined questions (e.g.,
diverticulosis, bowel-wall change, lesion extent, hernia site, acute-vs-chronic).
Prognosis and acquisition questions are excluded. By answer format the benchmark
contains 187{,}189 yes/no, 79{,}711 multiple-choice, and 13{,}211 severity items.

\begin{table}[t]
\centering
\caption{Question families in the IID benchmark.}
\label{tab:families}
\begin{tabular}{llrl}
\toprule
Family & Answer type & Count & Description \\
\midrule
Finding   & yes/no        & 174{,}093 & presence of a specific abnormality \\
Organ     & multiple choice & 60{,}014 & most prominent organ abnormality \\
Extended  & yes/no + choice & 26{,}553 & report-mined auxiliary questions \\
Score     & ordinal (L0--L4) & 13{,}211 & severity grading \\
Triage    & multiple choice & 4{,}997 & acute intra-abdominal process \\
Side      & multiple choice & 1{,}243 & laterality \\
\midrule
Total     &               & 280{,}111 & \\
\bottomrule
\end{tabular}
\end{table}

\subsection{Splits and leakage control}
\label{sec:splits}
Splits are assigned at the \textbf{study level} and inherited from the source
cohort's partition. We verify on the released files that the training/validation
pool and the test set share \textbf{no study and no image volume} (study overlap
$=0$, image overlap $=0$), so no report or scan is seen in more than one split.
Table~\ref{tab:splits} gives the per-split composition.
% NOTE (honesty): patient identifiers are not available, so patient-level grouping
% cannot be enforced; the source flags patient_leakage_possible=True. The claim above
% is strictly study/image-level. Do NOT upgrade this to "patient-level" without a
% patient-id dedup pass.

\begin{table}[t]
\centering
\caption{Per-split composition of the IID benchmark.}
\label{tab:splits}
\begin{tabular}{lrrrrr}
\toprule
Split & Studies & Questions & Yes/No & Choice & Score \\
\midrule
Train & 15{,}178 & 168{,}093 & 112{,}322 & 47{,}798 & 7{,}973 \\
Val   & 5{,}009  & 55{,}175  & 36{,}926  & 15{,}702 & 2{,}547 \\
Test  & 5{,}080  & 56{,}825  & 37{,}925  & 16{,}210 & 2{,}690 \\
\midrule
Total & 25{,}270 & 280{,}111 & 187{,}189 & 79{,}711 & 13{,}211 \\
\bottomrule
\end{tabular}
\end{table}

\subsection{Organ and disease coverage}
Organ questions cover eight organ systems (Table~\ref{tab:organs}); four systems
whose option set degenerates to a single abnormality (spleen, pancreas, lymph
nodes, prostate) are removed at release. Finding questions cover \textbf{30
abnormalities}; Table~\ref{tab:diseases} lists their positive/negative counts.
Overall the finding family is 37.4\% positive (65{,}184) and 62.6\% negative
(108{,}893). In the retained organ set, \texttt{no\_abnormality} accounts for
$\sim$81\% of items, so we down-sample it at release to avoid a trivial
majority-class baseline.

\begin{table}[t]
\centering
\caption{Organ-system distribution (organ family).}
\label{tab:organs}
\begin{tabular}{lr}
\toprule
Organ system & Count \\
\midrule
Liver \& biliary tree      & 12{,}683 \\
Lower thorax               & 9{,}543 \\
Gastrointestinal tract     & 8{,}140 \\
Kidneys \& ureters         & 7{,}842 \\
Musculoskeletal            & 6{,}981 \\
Peritoneal cavity          & 6{,}398 \\
Vasculature                & 5{,}486 \\
Gallbladder                & 2{,}940 \\
\bottomrule
\end{tabular}
\end{table}

\begin{table}[t]
\centering
\caption{Finding-family distribution over 30 abnormalities (positive / negative).}
\label{tab:diseases}
\begin{tabular}{lrrr}
\toprule
Abnormality & Total & Pos & Neg \\
\midrule
Gallstones                    & 15{,}085 & 1{,}985 & 13{,}100 \\
Atelectasis                   & 12{,}271 & 11{,}720 & 551 \\
Renal hypodensities           & 11{,}929 & 5{,}630 & 6{,}299 \\
Surgically absent gallbladder & 11{,}689 & 577 & 11{,}112 \\
Bowel obstruction             & 11{,}297 & 564 & 10{,}733 \\
Hydronephrosis                & 10{,}615 & 935 & 9{,}680 \\
Biliary ductal dilation       & 10{,}040 & 2{,}625 & 7{,}415 \\
Renal cyst                    & 9{,}247 & 8{,}086 & 1{,}161 \\
Metastatic disease            & 8{,}659 & 792 & 7{,}867 \\
Anasarca                      & 8{,}409 & 1{,}465 & 6{,}944 \\
Pleural effusion              & 8{,}160 & 3{,}809 & 4{,}351 \\
Atherosclerosis               & 7{,}656 & 7{,}264 & 392 \\
Free air                      & 5{,}997 & 579 & 5{,}418 \\
Abdominal aortic aneurysm     & 5{,}784 & 219 & 5{,}565 \\
Hepatic steatosis             & 4{,}898 & 2{,}646 & 2{,}252 \\
Lymphadenopathy               & 4{,}856 & 979 & 3{,}877 \\
Ascites                       & 3{,}476 & 2{,}577 & 899 \\
Fracture                      & 3{,}194 & 1{,}345 & 1{,}849 \\
Submucosal edema              & 3{,}126 & 1{,}667 & 1{,}459 \\
Appendicitis                  & 2{,}956 & 133 & 2{,}823 \\
Pancreatic atrophy            & 2{,}933 & 1{,}695 & 1{,}238 \\
Cardiomegaly                  & 1{,}868 & 792 & 1{,}076 \\
Coronary calcification        & 1{,}776 & 1{,}644 & 132 \\
Hiatal hernia                 & 1{,}690 & 1{,}650 & 40 \\
Osteopenia                    & 1{,}627 & 1{,}023 & 604 \\
Splenomegaly                  & 1{,}601 & 876 & 725 \\
Hepatomegaly                  & 1{,}467 & 425 & 1{,}042 \\
Thrombosis                    & 869 & 707 & 162 \\
Prostatomegaly                & 659 & 650 & 9 \\
Aortic valve calcification    & 243 & 125 & 118 \\
\midrule
Total                         & 174{,}077 & 65{,}184 & 108{,}893 \\
\bottomrule
\end{tabular}
\end{table}

\subsection{Label verification}
\label{sec:verification}
Each candidate label is verified by a \textbf{two-plus-one LLM reader ensemble}.
Two base readers (Qwen3-32B and Phi-4) independently judge each deduplicated
report text unit, blind to the rule-derived label. A third reader
(Qwen3.8-27B) adjudicates only the tied cases that receive exactly one support
vote. A question is retained only if it obtains \textbf{at least two support votes
out of three} (a 2-of-3 agreement gate); organ questions are judged by their own
three-reader vote rather than borrowing the finding verdict.

Implied negatives (``the report does not mention the finding, therefore it is
absent'') are \emph{also} reviewed by all three readers. Because a direct yes/no
prompt would be leading, they are re-queried with neutral, non-leading templates
(neutral-template binding) and judged by the same ensemble under the same 2-of-3
support gate; on top of that gate, each binding must keep its reader opposition
rate $\le 20\%$. Recovered items are tagged distinctly
(\texttt{implied\_negative\_v2}). Labels backed solely by a normal organ section
(\texttt{organ\_backed\_negative}) are tagged separately and excluded from the
released finding set.

\paragraph{Filtering yield.}
The factory produces 733{,}452 candidate questions; the verification gate retains
395{,}971, i.e., \textbf{46.0\% are dropped}. The main rejection reasons are
boilerplate-only implied negatives (79{,}551), non-auditable items (73{,}519),
bindings exceeding the opposition threshold (53{,}775), opposed bindings
(48{,}126), and reader disagreement (44{,}742). Release assembly then removes
organ-backed negatives and degenerate organ systems and down-samples
\texttt{no\_abnormality}, yielding the final 280{,}111 items.

\subsection{Human audit}
We manually audit a random sample of \textbf{400 retained IID items}
(seed-fixed). Using the criterion that a label is correct if the evidence supports
it (or it is the best available multiple-choice option with reader agreement), we
find \textbf{0 incorrect labels}; 12 items (3.0\%) are flagged as borderline, of
which 7 reflect organ option-set coverage gaps, 3 mixed/ambiguous evidence, and 2
borderline reader disagreement. For the out-of-distribution set we audit
\textbf{200 items}, finding 4 incorrect and 8 borderline.

\subsection{Out-of-distribution evaluation}
The OOD benchmark is the the external OOD cohort abdominal CT cohort, held out entirely from
training and used for prediction only (no OOD label enters fitting and no OOD
readout drives model selection). It comprises \textbf{421 studies and 2{,}807
questions} (1{,}947 yes/no, 801 choice, 59 score) over 24 abnormalities.
Table~\ref{tab:oodiid} summarizes the distribution shift relative to IID.
% NOTE: pin the OOD total before submission (current frozen test = 2,807; earlier
% versions 2,577 / 2,823 exist).

\begin{table}[t]
\centering
\caption{IID vs.\ OOD distribution comparison.}
\label{tab:oodiid}
\begin{tabular}{lll}
\toprule
 & IID (the source cohort) & OOD (the external OOD cohort) \\
\midrule
Source            & the source abdominal CT cohort & the external OOD cohort (external) \\
Splits            & train / val / test & test only (prediction-only) \\
Studies           & 25{,}270 & 421 \\
Questions         & 280{,}111 & 2{,}807 \\
Abnormalities     & 30 & 24 \\
Organ systems     & 8 & 8 (incl.\ merged groups) \\
Answer types      & 187{,}189 / 79{,}711 / 13{,}211 & 1{,}947 / 801 / 59 \\
\bottomrule
\end{tabular}
\end{table}

The OOD set omits six IID abnormalities---five lower-thorax findings
(atelectasis, pleural effusion, coronary calcification, aortic valve
calcification, cardiomegaly) and anasarca---reflecting the external OOD cohort's abdominal focus.
Disease prevalence also differs markedly: IID is dominated by common/incidental
findings (gallstones, atelectasis, renal hypodensities), whereas OOD is enriched
for pathology (lymphadenopathy, ascites, splenomegaly, renal cyst, biliary ductal
dilation). OOD further uses merged organ groups with whole-image fallback where a
dedicated organ category is unavailable.